% appendix_beliefs.tex: the beliefs the estimates hold at their measured values, sized in the 2022 SCF (moved from Section 8 on 2026-09-16).
% Results: scf_beliefs_by_group.csv, scf_coholding_by_group.csv, scf_belief_corrections.csv (Code/scf/scf_beliefs.R); tables in Appendix appendix:scf.
\section{Beliefs: What the Estimates Hold Fixed}\label{appendix:beliefs}

The estimates of Sections \ref{sec:rate_people} and \ref{sec:quasi} condition on beliefs, and this appendix sizes what the beliefs of constrained households do to them. Proposition \ref{prop:beliefs} states the dependence: the weight on a payment due at $s$ contains the borrower's own probability of still facing it, $\tilde Q_s$, and the estimates substitute the objective survival $Q_s$ measured in the panel. Under constant hazards, $\tilde Q_s = e^{-\tilde h (s-1)/12}$ and $Q_s = e^{-h(s-1)/12}$ in equation \eqref{eq:window}, and the substitution is additive in the rate:
\begin{equation}\label{eq:beliefrate}
\rho_{\text{meas}} \;=\; \rho + \tilde h - h ,
\end{equation}
where $h$ is the annual hazard of leaving the loan and $\tilde h$ the borrower's subjective one; a borrower who expects to stay in the loan longer than the hazard implies ($\tilde h < h$) reveals a lower rate than she has, and one who expects to leave sooner reveals a higher one. The second belief enters the premium. Appendix \ref{appendix:mucorrection} makes $\lambda$ a function of the persistence of the distress state, and a borrower who expects her low income to be transitory weighs a later dollar less than the objective persistence implies. Under constant relative risk aversion $\gamma$ with consumption moving with income, the first-year part of her weight is
\begin{equation}\label{eq:beliefpremium}
\log \tilde\lambda \;=\; -\gamma\, \mathbb{E}^{i}[\Delta \log y] + \tfrac{1}{2}\gamma^{2}\, \mathrm{Var}^{i}[\Delta \log y] ,
\end{equation}
so an expected recovery of $g$ lowers the first-year weight by the factor $e^{-\gamma g}$ relative to a borrower who expects no change. By Theorem \ref{thm:betalambda} this part sits inside the measured product $\beta\lambda$, and no timing comparison separates it from liquidity or present bias. Neither $\tilde h$ nor $\mathbb{E}^{i}[\Delta \log y]$ is observed in the bureau's data. The Survey of Consumer Finances asks households what they expect of their own income, and the paper uses the 2022 survey to size both objects; Appendix \ref{appendix:scf} lists the questions, the construction of the groups, and the full tabulations (Tables \ref{tab:scf_beliefs} and \ref{tab:scf_coholding}).

Constrained households, those turned down for credit or fearing a turndown, sixty days late on a payment, or using a payday loan within the year, 21 percent of households, expect their income to recover: 29 percent report an unusually low year, against 14 percent of the unconstrained, and their mean expected log income change is $0.11$ (standard error $0.02$) against $0.02$ ($0.01$); the gradient in income is steeper, $0.22$ ($0.02$) in the bottom quintile and $-0.09$ ($0.01$) in the top. The same households are less certain of next year's income and budget over shorter horizons, and their expectations for the economy do not differ from the unconstrained. These are the beliefs the estimates hold at their objective values, and they differ across groups in the direction that would matter.

What the reported beliefs do to the estimates follows from the two equations. Through equation \eqref{eq:beliefrate}, optimism about one's own income is optimism about staying in the loan, and it moves the rate by at most the group's default hazard, the only part of $h$ that a borrower's own default can change: the annual default hazard in the auto panel is $0.06$ in the bottom score decile and $0.0002$ in the top, so the estimate for a bottom-decile borrower is within $0.06$ of her rate on this account whatever she believes about her own default, and the profile of Section \ref{sec:rate_people} across deciles is not a belief artifact on the default side. Through equation \eqref{eq:beliefpremium} with $\gamma = 1.543$, the value at which the calibration of $\lambda(p)$ passes through Indarte's premium, the reported expectation lowers the constrained group's first-year weight by the factor $e^{-1.543 \times 0.11} = 0.84$, a 16 percent first-year discount, and the bottom quintile's by $0.72$, a 28 percent discount, against 3 percent for the unconstrained. In log units the belief-driven part of the premium is $0.17$ for the constrained and $0.33$ for the bottom quintile, of a total that runs from $0.95$ to $1.39$ across the bureau's groups: a sizeable part of the front-loading of constrained and low-income borrowers is expected recovery rather than liquidity or impatience, and it is a property of the premium, not of the rate. On the exit side the correction is bounded by nothing the survey reports: a borrower who expects to refinance or trade in at thirty months when the measured hazard is $0.19$ has $\tilde h = 0.40$ and reveals a rate $0.21$ above her own, the size of the middle-decile estimates, and the survey asks for the expected repayment date only when a loan is off schedule. This is the one belief the paper cannot bound, and Section \ref{sec:rate_people} states it where the middle deciles are reported.

The survey's revealed-preference content is the comparison of Corollary \ref{cor:separate}, two dollars at one date rather than one dollar at two: among the 44 percent of households that revolve a card balance, 63 percent (standard error 1.1) hold liquid assets at least equal to the balance at a median card rate of 17 percent, and for them the ratio of the borrowing and saving first-order conditions cancels $\rho$ and bounds the liquidity-service value of a dollar in the account at about 1.4 percent per month (Table \ref{tab:scf_coholding}). The margin is present in every group and measures the liquidity part of the premium for the 28 percent of households on it; it says nothing about the rate.
